% from ega4-17-fr.tex
\oldpage[IV-4]{92}

(iv) La condition b) de (17.13.2) exprime encore que l'homomorphisme $\alpha:g^*(\sh{N}_{Y/X})\to\Omega^1_{X'/S}\otimes_{\sh{O}_{X'}}\sh{O}_{Y'}$ est \emph{universellement injectif} relativement à $Y'$ (II, 9.18).
\end{remarks}

\begin{env}[17.13.5]
\label{IV.17.13.5}
Considérons maintenant un préschéma $S$, trois $S$-préschémas $X$, $Y$, $Z$, deux $S$-morphismes $f:Y\to X$, $g:Z\to X$; posons $T=Y\times_XZ$; on sait alors (I, 5.3.5) qu'on a un diagramme commutatif
\[
\label{IV.17.13.5.1}
\tag{17.13.5.1}
\xymatrix{
  X \ar[d]^{\Delta} & T=Y\times_X Z \ar[l] \ar[d]^{j}\\
  X\times_S X & Y\times_S Z \ar[l]^{u}
}
\]
faisant de $T$ le produit des $(X\times_SX)$-préschémas $X$ et $Y\times_SZ$, où $u=f\times_Sg$. Comme $\Delta$ est une $S$-immersion (I, 5.3.9), on est dans la situation du diagramme (17.13.1.1); ce qui correspond à $\Omega^1_{X'/S}$ dans (17.13.1) est alors $(\Omega^1_{Y/S}\otimes_{\sh{O}_Y}\sh{O}_{Y\times_SZ})\oplus(\Omega^1_{Z/S}\otimes_{\sh{O}_Z}\sh{O}_{Y\times_SZ})$ en vertu de (16.4.23). D'autre part, ce qui correspond au $\sh{O}_Y$-Module $\sh{N}_{Y/X}$ dans (17.13.1) est ici par définition $\Omega^1_{X/S}$ (16.3.1); il correspond donc à (17.13.1.3) une suite exacte
\[
\label{IV.17.13.5.2}
\tag{17.13.5.2}
\Omega^1_{X/S}\otimes_{\sh{O}_X}\sh{O}_T\xrightarrow{\rho}
(\Omega^1_{Y/S}\otimes_{\sh{O}_Y}\sh{O}_T)\oplus
(\Omega^1_{Z/S}\otimes_{\sh{O}_Z}\sh{O}_T)
\xrightarrow{\sigma}\Omega^1_{T/S}\to0
\]
où il reste à préciser les homomorphismes $\rho$ et $\sigma$. Tenant compte d'abord de (16.4.23) et de (0, 20.5.2), on voit que si $p:T\to Y$, $q:T\to Z$ sont les projections canoniques, on a, avec les notations de (16.4.19),
\[
\label{IV.17.13.5.3}
\tag{17.13.5.3}
\sigma=p_{T/Y/S}+q_{T/Z/S}.
\]
D'autre part, pour évaluer $\rho$, utilisons la commutativité du carré de gauche dans (17.13.1.2), ce qui, dans le cas présent, ramène d'abord à expliciter l'homomorphisme canonique
\[
\rho':\Omega^1_{X/S}\to\Omega^1_{(X\times_SX)/S}\otimes_{\sh{O}_{X\times_SX}}\sh{O}_X
\]
défini dans (16.4.21) appliqué à l'immersion $\Delta$. On peut se borner au cas où $S=\Spec(A)$, $X=\Spec(B)$ sont affines, et alors $\rho'$ correspond à l'homomorphisme $\delta$ de (0, 20.5.11.2), où il faut remplacer $B$ par $B\otimes_AB$ et $C$ par $B=(B\otimes_AB)/\mathfrak J_{B/A}$. On voit alors que $\delta$ fait correspondre à la classe de $x\otimes1-1\otimes x\pmod{\mathfrak J_{B/A}^2}$ (pour un $x\in B$) l'image de
\[
(x\otimes1-1\otimes x)\otimes(1\otimes1)-(1\otimes1)\otimes(x\otimes1-1\otimes x)
\]
dans $\Omega^1_{(B\otimes_AB)/A}\otimes_{B\otimes_AB}B$; mais l'élément précédent peut s'écrire
\[
\begin{split}
&((x\otimes1)\otimes(1\otimes1)-(1\otimes1)\otimes(x\otimes1))\\
&\quad-((1\otimes x)\otimes(1\otimes1)-(1\otimes1)\otimes(1\otimes x))
\end{split}
\]
et on voit donc que $\rho'$ est la différence des deux homomorphismes $\pi^{(1)}$ et $\pi^{(2)}$ de $\Omega^1_{X/S}$ dans $\Omega^1_{(X\times_SX)/S}\otimes_{\sh{O}_{X\times_SX}}\sh{O}_X$, correspondant respectivement à la première et à la seconde projection de $X\times_SX$ dans $X$, par (16.4.3.3). Pour obtenir $\rho$, il faut d'abord considérer l'homomorphisme
\[
\rho'':\Omega^1_{(X\times_SX)/S}\otimes_{\sh{O}_{X\times_SX}}\sh{O}_{Y\times_SZ}\to\Omega^1_{(Y\times_SZ)/S}
\]
correspondant par (16.4.3.3) au morphisme $u$ de (17.13.5.1), puis, après tensorisation par $\sh{O}_T$, former le composé $(\rho''\otimes1)\circ(\rho'\otimes1)$; il résulte de ce qui précède que l'on a, avec les notations de (16.4.18),
\[
\label{IV.17.13.5.4}
\tag{17.13.5.4}
\rho=(f_{Y/X/S}\otimes1_{\sh{O}_T})\oplus(-g_{Z/X/S}\otimes1_{\sh{O}_T}).
\]

